\documentclass[11pt,a4paper]{article}

% ---------------------------------------------------------------
% Packages
% ---------------------------------------------------------------
\usepackage[T1]{fontenc}
\usepackage{amsmath, amsthm, mathtools}
\usepackage{libertinus}
\usepackage{libertinust1math}
\usepackage{bm}
\usepackage{graphicx}
\usepackage{xcolor}
\usepackage{tikz}
\usetikzlibrary{arrows.meta, patterns}
\usepackage[most]{tcolorbox}
\usepackage{siunitx}
\sisetup{group-separator={\,}, group-minimum-digits=4, exponent-product=\cdot}
\DeclareSIUnit\angstrom{\text{\AA}}
\DeclareSIUnit\bohr{\ensuremath{a_0}}
\usepackage{booktabs}
\usepackage{tabularx}
\usepackage[numbers,sort&compress]{natbib}
\usepackage{microtype}
\usepackage{parskip}
\usepackage[margin=2.5cm, headheight=20pt]{geometry}
\usepackage{titlesec}
\usepackage{fancyhdr}
\usepackage{lastpage}
\usepackage{hyperref}
\usepackage[capitalise,noabbrev]{cleveref}

\numberwithin{equation}{section}

% ---------------------------------------------------------------
% Corporate colours (sampled from the FAU logo) and layout
% ---------------------------------------------------------------
\definecolor{faublue}{RGB}{0,47,108}
\definecolor{fautech}{RGB}{119,159,181}

\hypersetup{colorlinks=true, linkcolor=faublue, citecolor=faublue, urlcolor=faublue,
  pdftitle={Discrete Dislocation Dynamics},
  pdfauthor={Ali Ghelichkhani}}

\titleformat{\section}{\sffamily\Large\bfseries\color{faublue}}{\thesection}{0.75em}{}
\titleformat{\subsection}{\sffamily\large\bfseries\color{faublue}}{\thesubsection}{0.7em}{}
\titleformat{\paragraph}[runin]{\sffamily\bfseries\color{faublue}}{}{0pt}{}
\titlespacing*{\paragraph}{0pt}{1.2ex plus 0.4ex}{0.6em}

\pagestyle{fancy}
\fancyhf{}
\fancyhead[L]{\sffamily\small\color{fautech}Materials Simulation Practical~3 \textbullet{} Part~I}
\fancyhead[R]{\raisebox{-0.2\height}{\includegraphics[height=0.42cm]{figures/ww8.png}}}
\fancyfoot[L]{\sffamily\small\color{fautech}A.~Ghelichkhani}
\fancyfoot[R]{\sffamily\small\color{fautech}\thepage\ / \pageref*{LastPage}}
\renewcommand{\headrule}{\hbox to\headwidth{\color{fautech}\leaders\hrule height 0.4pt\hfill}}

\newtcolorbox{task}[1]{enhanced, sharp corners, frame hidden,
  colback=fautech!12, borderline west={2.5pt}{0pt}{faublue},
  left=9pt, right=9pt, top=6pt, bottom=6pt, after skip=12pt,
  before={\par\nopagebreak\vskip 6pt},
  fontupper=\sffamily\small,
  before upper={{\bfseries\color{faublue}Task~#1.}\enspace}}
\setcounter{MaxMatrixCols}{20}

\theoremstyle{remark}
\newtheorem{remark}{Remark}[section]

% ---------------------------------------------------------------
% Notation
% ---------------------------------------------------------------
\newcommand{\bigO}[1]{\mathcal{O}\!\left(#1\right)}
\newcommand{\tauc}{\tau_{\mathrm{c}}}
\newcommand{\bb}{\bm{b}}
\newcommand{\bl}{\bm{l}}

\begin{document}

\begin{titlepage}
  \noindent
  \includegraphics[height=2.0cm]{figures/fau-techfak.png}\hfill
  \raisebox{0.6cm}{\includegraphics[height=1.1cm]{figures/ww8.png}}

  \vspace{0.2cm}
  {\color{fautech}\rule{\textwidth}{0.6pt}}\\[0.1cm]
  {\sffamily\small\color{fautech}Institute of Materials Simulation (WW8)\hfill
    Winter semester 2026/27}

  \vspace*{3.4cm}
  {\sffamily\large\color{fautech}\textls[120]{MATERIALS SIMULATION PRACTICAL 3}\par}
  \vspace{0.5cm}
  {\sffamily\bfseries\color{faublue}\fontsize{30}{36}\selectfont
    Discrete Dislocation\\[0.1em] Dynamics\par}
  \vspace{0.5cm}
  {\sffamily\Large\color{faublue}Activation of a Frank--Read Source in Copper with microMegas\par}
  \vspace{0.7cm}
  {\color{faublue}\rule{3cm}{2pt}}\par
  \vspace{0.4cm}
  {\sffamily\large Report, Part~I \textbullet{} Theory and Results\par}

  \vfill
  {\sffamily
  \renewcommand{\arraystretch}{1.35}
  \begin{tabular}{@{}>{\small\color{fautech}}l@{\hspace{1.2em}}l@{}}
    Author     & Ali Ghelichkhani \\
    E-mail     & \href{mailto:ali.ghelichkhani@fau.de}{ali.ghelichkhani@fau.de} \\
    Supervisor & Leon Pyka \\
    Module     & Foundations of Materials Simulation \\
    Institute  & Institute of Materials Simulation (WW8), FAU Erlangen-N\"urnberg \\
    Part~II    & Simulations, Jupyter notebook \texttt{ddd\_frank\_read.ipynb}
  \end{tabular}\par}
\end{titlepage}

\section*{Abstract}
  We use the discrete dislocation dynamics code microMegas to determine the critical resolved shear
  stress $\tauc$ of a Frank--Read source on the $(\bar1\bar1\bar1)[10\bar1]$ slip system of a copper
  single crystal, and compare it with the line-tension estimate $\tauc = \alpha\mu b/L$. A
  $\SI{100}{\nano\metre}$ edge source operates at $\tauc = \SI{104.3}{\mega\pascal}$, i.e.\
  $\alpha = 0.99$, about $2.7$ times the large-bow-out estimate of Hirth and Lothe with a core radius of
  one Burgers vector. The value is insensitive to the number of time steps, to the box size between
  $1.5$ and $\SI{6}{\micro\metre}$ and, for this length, to the re-discretization length; it is
  dominated by the line-tension model, which changes it by a factor of two (Friedel:
  $\SI{46.7}{\mega\pascal}$, Foreman: $\SI{100.3}{\mega\pascal}$, DeWit: $\SI{104.3}{\mega\pascal}$). Over
  $L = 48$--$\SI{203}{\nano\metre}$ the critical stress falls as $L^{-0.72}$, weaker than the $1/L$ of a
  constant $\alpha$ and close to the $\ln(L/\rho)/L$ of the Hirth formula, so $\alpha$ increases with
  $L$ from $0.95$ to $1.38$. For fixed $L = \SI{100}{\nano\metre}$, $\alpha$ rises from $0.99$ for an edge
  source to $1.80$ for a screw source, reproducing the qualitative trend of orientation-dependent line
  tension. The scripts and the full set of runs are in Part~II.

\tableofcontents
\clearpage

% ===============================================================
\section{Introduction}\label{sec:intro}
% ===============================================================

Between the atomistic scale of molecular dynamics and the continuum scale of crystal plasticity lies
a range, roughly $\SI{100}{\nano\metre}$ to some tens of micrometres, in which the plastic response of a
crystal is governed by the motion and interaction of individual dislocations. Molecular dynamics cannot
reach these sizes and times, and continuum models do not resolve the dislocations whose collective
behaviour they average. Discrete dislocation dynamics (DDD) closes the gap~\citep{script3, fivel2008}:
it represents the dislocations as lines made of straight segments embedded in a linear elastic
continuum and moves them under the forces of the applied stress and of each other.

This report studies the simplest multiplication mechanism, the Frank--Read source. We review the DDD
algorithm and the microMegas code (\cref{sec:ddd}), derive the critical stress of a Frank--Read source
from the line-tension approximation (\cref{sec:fr}), and present the simulations: the set-up and the
determination of $\tauc$ and $\alpha$ for a $\SI{100}{\nano\metre}$ edge source (\cref{sec:task1}),
the dependence on the source length and on the simulation parameters (\cref{sec:task2}), and the
dependence on the character angle (\cref{sec:task3}).

\begin{table}[ht]
  \centering
  \caption{Correspondence between the tasks of the tutorial script and this report.}
  \label{tab:taskmap}
  \begin{tabularx}{0.9\linewidth}{lXl}
    \toprule
    Task & Content & Section \\
    \midrule
    1.1--1.5 & microMegas, wrapper, Cu crystal, $\SI{100}{\nano\metre}$ edge source & \ref{sec:setup} \\
    1.6 & estimate of the critical stress & \ref{sec:estimate} \\
    1.7, 1.8 & determination of $\tauc$; fit of $\alpha$ & \ref{sec:determination} \\
    2 & $\tauc$ vs.\ segment length; influence of simulation parameters & \ref{sec:task2} \\
    3 & $\tauc$ and $\alpha$ vs.\ character angle $\beta$ & \ref{sec:task3} \\
    \bottomrule
  \end{tabularx}
\end{table}

% ===============================================================
\section{Discrete dislocation dynamics}\label{sec:ddd}
% ===============================================================

Every DDD code repeats the same three steps~\citep{devincre2011}:
\begin{enumerate}\setlength{\itemsep}{1pt}
  \item evaluate the stress at every segment, the sum of the applied stress and the internal stress
    field of all other segments;
  \item evaluate the force per unit length on each segment from the Peach--Koehler formula,
    \begin{equation}\label{eq:pk}
      F = \left(\tau_{\mathrm{ext}} + \tau_{\mathrm{int}}\right)b\,l,
    \end{equation}
    in which only the glide component of the resolved shear stress enters for glide motion;
  \item move the segments with a mobility law $v = M(F)$ and an explicit Euler step,
    $r(t+\Delta t) = r(t) + v(t)\,\Delta t$, and update the topology: split segments that have grown
    too long, merge short ones, and handle junctions, annihilation and cross-slip.
\end{enumerate}
The internal stress is long-ranged ($\propto 1/r$), which makes step~1 the expensive one. The
self-interaction of a curved line is the delicate part: the stress of a segment on itself diverges and
must be regularized, either by a core cut-off in the full elastic interaction or, in the simpler
approach used for the local part of the self-force, by a \emph{line tension} that pulls a curved line
straight.

\paragraph{microMegas.}
microMegas (mM)~\citep{devincre2011} discretizes the lines on a lattice: each segment is a multiple of
one of eight elementary vectors per slip system --- pure screw, pure edge and two mixed orientations in
both senses --- so that a curved line is a polygon of these directions. The elementary screw vector is
\texttt{echelle} Burgers vectors long ($13.5\,b = \SI{3.45}{\nano\metre}$), the edge vector
$\SI{5.97}{\nano\metre}$. Segments that exceed a length \texttt{ldis\_act} are re-discretized. The local
line tension is chosen with the key \texttt{linten}: Friedel's constant $T = \mu b^2/2$, the DeWit--Koehler
form that depends on the character of the segment and on the logarithm of the local curvature radius,
or Foreman's variant. The mobility is viscous, $v = \tau b/B$, with a small friction stress of
$\SI{0.5}{\mega\pascal}$. The course provides a Python wrapper whose dataclass \texttt{SimulationConfig}
sets the segment length, the character angle $\beta$, the box size and the start point, and whose
nested \texttt{ContCuConfig} holds all 41 control parameters of the \texttt{ContCu} file; the wrapper
writes the three input files and launches the engine.

% ===============================================================
\section{The Frank--Read source}\label{sec:fr}
% ===============================================================

A segment of length $L$ pinned at both ends bows out under a resolved shear stress $\tau$. The
Peach--Koehler force per unit length, $\tau b$, acts normal to the line and is balanced by the line
tension $T$ of the curved arc, $\tau b = T/r$ for an arc of radius $r$. As $\tau$ grows, $r$ decreases
until the arc becomes a semicircle, $r = L/2$, the smallest radius a segment of length $L$ can have.
Beyond this point the equilibrium is unstable: the arc expands at decreasing curvature, its two arms
wrap around the pinning points, meet and annihilate, and a closed loop is released while the original
segment is restored. The source then repeats indefinitely. The critical stress is therefore
\begin{equation}\label{eq:taucsimple}
  \tauc = \frac{2T}{bL} = \alpha\,\frac{\mu b}{L},
\end{equation}
with $\alpha = 1$ for Friedel's line tension $T = \mu b^2/2$. A better line tension accounts for the
character of the line and for the logarithmic dependence of the elastic energy on its extent. Hirth and
Lothe~\citep{hirth1982} give for the large bow-out
\begin{equation}\label{eq:hirth}
  b\tauc = \frac{\mu b^2}{4\pi(1-\nu)\,r}
  \left\{\left[1 - \frac{\nu}{2}\left(3 - 4\cos^2\beta\right)\right]\ln\frac{L}{\rho} - 1 + \frac{\nu}{2}\right\},
  \qquad r = \frac{L}{2},
\end{equation}
where $\beta$ is the angle between the initial line and the Burgers vector and $\rho$ the inner cut-off
radius. Comparing with \cref{eq:taucsimple},
\begin{equation}\label{eq:alpha}
  \alpha(\beta, L) = \frac{1}{2\pi(1-\nu)}
  \left\{\left[1 - \frac{\nu}{2}\left(3 - 4\cos^2\beta\right)\right]\ln\frac{L}{\rho} - 1 + \frac{\nu}{2}\right\}.
\end{equation}
Two consequences matter for the simulations. First, $\alpha$ is not a constant but grows with
$\ln(L/\rho)$, so $\tauc$ decreases more slowly than $1/L$. Second, $\alpha$ is larger for a screw
source ($\cos\beta = 1$) than for an edge source ($\cos\beta = 0$): the arms of a bowing edge source are
screw-like, and screw segments have the higher line tension~\citep{anderson2017}. For $\nu = 0.33$ and
$L = 1000\rho$ the script quotes $\alpha_{\mathrm{edge}} \approx 0.5$ and
$\alpha_{\mathrm{screw}} \approx 1.5$; evaluating \cref{eq:alpha} directly gives $0.63$ and $1.71$.

% ===============================================================
\section{The $\SI{100}{\nano\metre}$ edge source (Task 1)}\label{sec:task1}
% ===============================================================

\subsection{Set-up (Tasks 1.1--1.5)}\label{sec:setup}

\begin{task}{1}
  Clone the microMegas repository, read the instructions and identify the parameters of the
  \texttt{SimulationConfig} dataclass. Set up a Cu single crystal in a
  $\SI{3}{\micro\metre}\times\SI{3}{\micro\metre}\times\SI{3}{\micro\metre}$ box with a pinned
  $\SI{100}{\nano\metre}$ edge segment in the $(\bar1\bar1\bar1)$ plane with Burgers vector $[10\bar1]$.
  Estimate the critical activation stress, determine it from simulations with stresses around the
  estimate, and determine $\alpha$.
\end{task}

The reference slip system of the material file \texttt{dd/in/Cu} is exactly $(\bar1\bar1\bar1)[10\bar1]$,
so a segment generated with $\beta = 90^\circ$ is the requested edge source. Its length is a multiple of
the elementary edge vector: 17 of them give $L = \SI{101.5}{\nano\metre}$, which we use in all
evaluations. The material data are $\mu = \SI{42}{\giga\pascal}$, $\nu = 0.431$ and
$b = \SI{0.2553}{\nano\metre}$, so $\mu b = \SI{10.72}{\giga\pascal\nano\metre}$. With \texttt{shear = T}
the applied value \texttt{sigma0} is the resolved shear stress on the most highly stressed system; the
run log confirms this by printing \texttt{Tau = sigma0} and the uniaxial equivalent $\sigma_0/0.408$,
the Schmid factor of $\{111\}\langle110\rangle$ for a $[001]$ axis.

Running the wrapper on macOS required two fixes to the environment and two to the wrapper, all
documented in Part~II: the precompiled Linux object files in \texttt{dd/bin/mm.objs} had to be removed
before \texttt{make mm}; the output directories \texttt{dd/out/\{debug,vtk,SigInt\_dir,GammaboxVTK\}} had
to be created; \texttt{MicroMegasWrapper.run} passes the box size as a tuple that the segment generator
then modifies in place, which raises an exception; and the generator writes the box size as floating
point numbers, which the Fortran list-directed read rejects.

The simulations are run at constant stress (\texttt{mode\_deformation = 4}). This gives a sharp
criterion: below $\tauc$ the line bows out to a stable arc within a few hundred steps, the number of
segments stays at five to seven and the plastic strain saturates near $10^{-6}\,\%$; above $\tauc$ the
source operates, emits loops, the number of segments grows to hundreds and the plastic strain increases
without bound. The source is classified as active if the number of segments exceeds three times its
initial value plus thirty, and $\tauc$ is bracketed by a geometric bisection to $1\,\%$.

\subsection{Estimate (Task 1.6)}\label{sec:estimate}

With the microMegas values and \cref{eq:alpha} the estimates for $L = \SI{101.5}{\nano\metre}$ are
\begin{equation}\label{eq:estimate}
  \tauc \approx
  \begin{cases}
    \SI{39}{\mega\pascal} & (\alpha = 0.37,\ \rho = b,\ \nu = 0.431),\\
    \SI{49}{\mega\pascal} & (\alpha = 0.46,\ L = 1000\rho,\ \nu = 0.431),\\
    \SI{55}{\mega\pascal} & (\rho = b,\ \nu = 0.33),
  \end{cases}
\end{equation}
which is the 50--65\,MPa range quoted in the wrapper documentation. The simulations were started around
these values and the bracket widened upwards when the source did not operate.

\subsection{Determination of $\tauc$ and $\alpha$ (Tasks 1.7, 1.8)}\label{sec:determination}

The source does not operate at $\SI{60}{}$, $\SI{80}{}$ or $\SI{95}{\mega\pascal}$; the bisection gives
\begin{equation}\label{eq:result1}
  \boxed{\;\tauc = \SI{104.3}{\mega\pascal},\qquad
  \alpha = \frac{\tauc L}{\mu b} = 0.99\;}
\end{equation}
for the default control parameters, about $2.7$ times the estimate with $\rho = b$. To find the origin of
the difference, \cref{tab:robust} repeats the determination with one parameter changed at a time.

\begin{table}[ht]
  \centering
  \caption{Critical stress of the $\SI{101.5}{\nano\metre}$ edge source for variations of the control
    parameters (Part~II).}
  \label{tab:robust}
  \begin{tabular}{lrr}
    \toprule
    Variation & $\tauc$ [MPa] & $\alpha$ \\
    \midrule
    reference (DeWit line tension, 4000 steps, $\SI{3}{\micro\metre}$ box) & 104.3 & 0.99 \\
    12\,000 steps & 104.3 & 0.99 \\
    box $\SI{1.5}{\micro\metre}$ and $\SI{6}{\micro\metre}$ & 104.3 & 0.99 \\
    re-discretization length $\SI{0.1}{\micro\metre}$ & 104.3 & 0.99 \\
    Foreman line tension & 100.3 & 0.95 \\
    Friedel line tension & 46.7 & 0.44 \\
    \bottomrule
  \end{tabular}
\end{table}

Neither the run length nor the box size has any influence: a single source in a micrometre-sized box
does not feel its periodic images, and four thousand steps of $\SI{1.25}{\pico\second}$ are ample for a
line moving at the viscous velocity $\tau b/B$ of several hundred metres per second. What controls
$\tauc$ is the line tension. With Friedel's constant line tension the source operates at
$\SI{47}{\mega\pascal}$, within $20\,\%$ of the Hirth estimate, whereas the default DeWit--Koehler form and
Foreman's variant, which raise the line tension of the screw-like arms of the bowing edge source, double
the critical stress. In other words, the critical stress of a source in DDD is essentially an input of
the model --- it is set by how the self-interaction of the curved line is regularized --- and the value
$\alpha \approx 1$ is a property of microMegas' default line tension as much as of copper.

% ===============================================================
\section{Dependence on the segment length (Task 2)}\label{sec:task2}
% ===============================================================

\begin{task}{2}
  Starting from Task~1, vary the length of the initial segment from $\SI{50}{}$ to
  $\SI{200}{\nano\metre}$, determine $\tauc$ for each length, plot the results together with the
  analytical solution and the relative error, explain for which length simulation and theory agree
  best, and investigate which simulation parameters (box size, minimum segment length, \dots) affect the
  results.
\end{task}

\begin{table}[ht]
  \centering
  \caption{Critical stress against segment length for the edge source (Part~II). ``Eq.~(6)'' is
    $\alpha\mu b/L$ with $\alpha = 0.99$ from Task~1; ``Hirth'' is \cref{eq:hirth} with $\rho = b$.}
  \label{tab:length}
  \begin{tabular}{rrrrrrrr}
    \toprule
    $L$ [nm] & $\tauc$ DeWit & eq.~(6) & error & Hirth & error & $\alpha$ & $\tauc$ Friedel \\
    \midrule
     47.7 & 213.0 & 221.7 & $-3.9\,\%$  & 66.9 & $+218\,\%$ & 0.95 & 83.9 \\
     77.6 & 121.1 & 136.4 & $-11.3\,\%$ & 47.8 & $+153\,\%$ & 0.88 & 50.6 \\
    101.5 & 104.3 & 104.3 & $0.0\,\%$   & 39.3 & $+165\,\%$ & 0.99 & 46.7 \\
    125.3 &  91.7 &  84.5 & $+8.6\,\%$  & 33.6 & $+173\,\%$ & 1.07 & 45.8 \\
    149.2 &  83.1 &  70.9 & $+17.1\,\%$ & 29.5 & $+182\,\%$ & 1.16 & 45.8 \\
    173.1 &  77.5 &  61.2 & $+26.7\,\%$ & 26.3 & $+194\,\%$ & 1.25 & 45.8 \\
    202.9 &  73.0 &  52.2 & $+40.0\,\%$ & 23.3 & $+214\,\%$ & 1.38 & 45.8 \\
    \bottomrule
  \end{tabular}
\end{table}

\Cref{tab:length} summarizes the results; Part~II plots them with the two analytical curves and the
relative errors. With the default line tension $\tauc$ decreases monotonically from $\SI{213}{}$ to
$\SI{73}{\mega\pascal}$, but more slowly than $1/L$: a power-law fit gives $\tauc \propto L^{-0.72}$.

\paragraph{Where simulation and theory agree.}
Against \cref{eq:taucsimple} with the $\alpha$ fitted at $\SI{100}{\nano\metre}$ the agreement is by
construction exact at $\SI{100}{\nano\metre}$ and degrades in both directions, to $-4\,\%$ at
$\SI{48}{\nano\metre}$ and $+40\,\%$ at $\SI{203}{\nano\metre}$. The growth of the error with $L$ is
physical: a constant $\alpha$ misses the logarithm in \cref{eq:alpha}, which makes long sources
relatively harder to operate. The Hirth formula contains that logarithm, and indeed the simulated
exponent $-0.72$ is close to its effective exponent of about $-0.8$ over this range; its absolute level,
however, is a factor $2.5$--$3$ too low for the DeWit line tension, as discussed in \cref{sec:determination}.
With Friedel's line tension the simulation agrees best with Hirth at $L \approx \SI{78}{\nano\metre}$
($+6\,\%$) and $\SI{100}{\nano\metre}$ ($+19\,\%$). For short sources the agreement deteriorates because
the line is built from only eight to seventeen elementary vectors, so the semicircular critical
configuration is represented by a handful of segments and its curvature is poorly resolved. For sources
longer than about $\SI{125}{\nano\metre}$ the Friedel results stop decreasing and stay at
$\SI{46}{\mega\pascal}$ regardless of $L$; switching cross-slip off does not change this, and the friction
stress of the mobility law is only $\SI{0.5}{\mega\pascal}$, so we attribute this floor to the lattice
discretization of the line rather than to physics. The best overall agreement is therefore found for
intermediate lengths of $\SI{75}{}$--$\SI{100}{\nano\metre}$, long enough to resolve the bow-out and
short enough to stay clear of the floor.

\paragraph{Which parameters matter.}
Part~II repeats the determination for $L = \SI{50}{}$ and $\SI{200}{\nano\metre}$ with modified
parameters:
\begin{itemize}\setlength{\itemsep}{2pt}
  \item \emph{Box size} ($\SI{1.5}{\micro\metre}$ instead of $\SI{3}{\micro\metre}$): no effect at either
    length; the source is far smaller than the box.
  \item \emph{Re-discretization length} \texttt{ldis\_act} ($\SI{0.1}{\micro\metre}$ instead of
    $\SI{0.25}{\micro\metre}$): $-27\,\%$ at $\SI{50}{\nano\metre}$, $-4\,\%$ at $\SI{200}{\nano\metre}$.
    Finer re-discretization lets the short source bow into a better-resolved arc.
  \item \emph{Elementary segment length} \texttt{echelle} ($10\,b$ instead of $13.5\,b$): $+274\,\%$ at
    $\SI{50}{\nano\metre}$, $+13\,\%$ at $\SI{200}{\nano\metre}$. Changing the lattice spacing of the
    discretization changes the realized geometry and the local line tension.
  \item \emph{Line tension model}: a factor of two at all lengths (\cref{tab:robust,tab:length}).
  \item \emph{Number of steps}: no effect once the bowed configuration has relaxed.
\end{itemize}
The minimum segment length and the line-tension treatment are the parameters that matter; the
dependence on them is strongest for short sources, which are represented by the fewest segments.

% ===============================================================
\section{Dependence on the line orientation (Task 3)}\label{sec:task3}
% ===============================================================

\begin{task}{3}
  Starting from Task~1, change the initial line orientation gradually from edge to screw by changing
  the angle $\beta$ between line and Burgers vector, keeping the line on the slip plane and composing it
  of pure edge, screw or mixed segments. Determine $\tauc$ and $\alpha$ for each orientation and plot
  $\alpha$ against $\beta$.
\end{task}

For $\beta = 0^\circ$, $60^\circ$ and $90^\circ$ the line is a single native segment of the
corresponding character. For intermediate angles the wrapper's generator composes a ``staircase'' of
$N_e$ edge and $N_s$ screw elementary vectors with $N_e/N_s$ chosen from
$L\sin\beta/l_{\mathrm{edge}}$ and $L\cos\beta/l_{\mathrm{screw}}$, distributed evenly along the line
(a Bresenham sequence) and grouped into blocks. Because the edge and screw vectors are perpendicular in
the slip plane, the realized end-to-end length and angle are
$L = \sqrt{(N_e l_{\mathrm{edge}})^2 + (N_s l_{\mathrm{screw}})^2}$ and
$\beta = \arctan(N_e l_{\mathrm{edge}}/N_s l_{\mathrm{screw}})$, which differ from the targets by up to
one degree and two nanometres; $\alpha$ is evaluated with the realized values.

\begin{table}[ht]
  \centering
  \caption{Critical stress and geometric factor against character angle for $L \approx
    \SI{100}{\nano\metre}$ (Part~II). The Hirth values use \cref{eq:alpha} with $\rho = b$ and
    $\nu = 0.431$.}
  \label{tab:beta}
  \begin{tabular}{rrrrrr}
    \toprule
    $\beta$ target & $\beta$ realized & $L$ [nm] & $\tauc$ [MPa] & $\alpha$ mM & $\alpha$ Hirth \\
    \midrule
    90 & 90.0 & 101.5 & 104.3 & 0.99 & 0.37 \\
    75 & 73.9 &  99.4 & 121.1 & 1.12 & 0.48 \\
    60 & 60.0 &  99.9 & 101.3 & 0.94 & 0.73 \\
    45 & 44.7 & 101.8 & 133.7 & 1.27 & 1.10 \\
    30 & 29.0 &  98.5 & 132.4 & 1.22 & 1.47 \\
    15 & 13.9 &  99.4 & 155.1 & 1.44 & 1.73 \\
     0 &  0.0 &  99.9 & 192.9 & 1.80 & 1.81 \\
    \bottomrule
  \end{tabular}
\end{table}

\Cref{tab:beta} lists the results; Part~II plots $\alpha(\beta)$ with the Hirth curves. The trend is the
expected one: $\alpha$ grows from $0.99$ for the edge source to $1.80$ for the screw source, a ratio of
$1.8$. In the semicircular critical configuration the arc of an edge source has its screw-like
parts near the pinning points, whereas a screw source has them at the apex, where the curvature drives
the instability; with an orientation-dependent line tension, which is highest for screw character,
the screw source is therefore the harder one to activate~\citep{anderson2017}. The Hirth formula predicts
a much stronger dependence, a ratio of $4.9$ for $\rho = b$: the two agree almost exactly for the screw source ($1.80$
against $1.81$), but microMegas puts the edge source $2.7$ times higher. Two points deviate from a smooth curve. The
$60^\circ$ source, built from one native mixed segment, is easier to activate than the staircase sources
on either side of it, and the $75^\circ$ staircase is harder than the pure edge. Both reflect the
discretization: a native segment and a staircase of the same nominal angle are different initial
configurations, and the steps of a staircase act as additional curvature that the line has to remove
before it can bow out smoothly.

% ===============================================================
\section{Conclusion}\label{sec:conclusion}
% ===============================================================

A Frank--Read source in copper simulated with microMegas operates at a well-defined critical resolved
shear stress that can be bracketed to one percent by constant-stress runs. For a
$\SI{100}{\nano\metre}$ edge source on $(\bar1\bar1\bar1)[10\bar1]$ we find
$\tauc = \SI{104}{\mega\pascal}$ and $\alpha = 0.99$. The critical stress decreases as $L^{-0.72}$
between $50$ and $\SI{200}{\nano\metre}$, so $\alpha$ is not a constant but grows with the source length,
as the logarithm in the Hirth--Lothe line-tension formula requires; and $\alpha$ increases from edge to
screw sources by a factor of $1.8$. The absolute level, however, is set by the numerical model rather
than by the crystal: the line-tension treatment changes $\tauc$ by a factor of two, and the
elementary segment length and re-discretization length change it by tens of percent for short sources,
whereas the box size and the run length have no influence at all. The practical lesson is that
the critical stress of a source in DDD has to be calibrated --- against atomistic simulations or
experiments --- before the code is used to predict the strength of a dislocation microstructure.

\begin{thebibliography}{99}

\bibitem{script3}
Institute of Materials Simulation (WW8),
``Discrete dislocation dynamics simulation,''
Materials Simulation Practical, tutorial script,
Friedrich-Alexander-Universit\"at Erlangen-N\"urnberg.

\bibitem{fivel2008}
M.~Fivel,
``Discrete dislocation dynamics: Principles and recent applications,''
in \textit{Multiscale Modeling of Heterogenous Materials}, O.~Cazacu (Ed.),
Wiley, 2008, pp.~17--36.

\bibitem{devincre2011}
B.~Devincre, R.~Madec, G.~Monnet, S.~Queyreau, R.~Gatti and L.~Kubin,
``Modeling crystal plasticity with dislocation dynamics simulations: The `microMegas' code,''
in \textit{Mechanics of Nano-objects}, Presses des Mines, Paris, 2011, pp.~81--100.

\bibitem{hirth1982}
J.~P.~Hirth and J.~Lothe,
\textit{Theory of Dislocations},
2nd~ed., Wiley, New York, 1982.

\bibitem{anderson2017}
P.~M.~Anderson, J.~P.~Hirth and J.~Lothe,
\textit{Theory of Dislocations},
3rd~ed., Cambridge University Press, Cambridge, 2017.

\bibitem{hull2011}
D.~Hull and D.~J.~Bacon,
\textit{Introduction to Dislocations},
5th~ed., Butterworth-Heinemann, Oxford, 2011.

\end{thebibliography}

\end{document}
